\section{Conclusiones}

\begin{itemize}
    \item Se logró analizar el movimiento de dos cuerpos con diferentes tipos de trayectoria: el automóvil con Movimiento Rectilíneo Uniforme (MRU) y el camión con Movimiento Rectilíneo Uniformemente Acelerado (MRUA), evidenciando las diferencias fundamentales entre velocidad constante y aceleración constante.
    
    \item Se determinó que ambos móviles se encuentran en dos instantes distintos, siendo el primero el más relevante físicamente ($t \approx 2{,}3$ s), lo cual demuestra que en sistemas con aceleración pueden existir múltiples puntos de intersección.
    
    \item Se comprobó que, aunque en $t = 10$ s las velocidades de ambos cuerpos son iguales, sus posiciones son diferentes, lo que confirma que igualdad de velocidad no implica igualdad de posición.
    
    \item El uso de herramientas matemáticas como derivadas e integrales permitió obtener de forma rigurosa las ecuaciones de velocidad y posición, facilitando la resolución del problema.
    
    \item La simulación en GeoGebra permitió validar gráficamente los resultados analíticos, reforzando la comprensión del comportamiento del movimiento y la interpretación física de las gráficas.
\end{itemize}